\chapter{A Concrete Naming Certificate for $\InterHistRateBoundaryUp$}
\label{ch:concrete-instances-inter-hist-rate-boundary-namecert}
\origin{ai}

The $\InterHistRateBoundaryUp$ packet realizes the inter-Hist maximum-rate boundary of \autoref{ch:visions-inter-hist-locality} as finite BEDC data. It binds the cross-history rate surface of \autoref{ch:concrete-instances-crosshistcausalrate-namecert}, the read gate of \autoref{ch:concrete-instances-maxratereadgate-namecert}, and the refusal surface of \autoref{ch:concrete-instances-noglobalsyncboundary-namecert} into one local naming carrier, so downstream consumers receive a gated route rather than a host speed, metric row, clock, or global frame.

\falsifiablePrediction{An accepted rate-boundary consumer that reads a maximum-rate row without both observer-symmetry and no-global-frame refusal rows refutes the packet.}

\independenceWitness{crosshistcausalrate, maxratereadgate, noglobalsyncboundary, interhistinvariant, crosshistcausalroute: the boundary packet is not bijective to any one sibling because it carries the shared gate plus route plus refusal surface.}

\begin{definition}[InterHistRateBoundary carrier]
\label{def:inter-hist-rate-boundary-carrier}
An $\InterHistRateBoundaryUp$ carrier is a finite $\BHist$ packet
$$
B=(A,D,G,R,S,F,I,H,Q,C,P,N)
$$
with the following visible rows:
\begin{enumerate}
\item $A$ is the paired observer-history boundary row that supplies the two displayed inter-Hist sides;
\item $D$ is the causal-route row, read through $\CrossHistCausalRouteUp$ and $\CrossHistCausalRateUp$ rather than through an ambient spacetime path;
\item $G$ is the accepted $\MaxRateReadGateUp$ row binding the rate read to the gate surface of \autoref{ch:concrete-instances-maxratereadgate-namecert};
\item $R$ is the maximum-rate row tied to $\MaxCausalRateUp$ as in \autoref{ch:concrete-instances-maxcausalrate-namecert};
\item $S$ is the observer-symmetry row for the displayed histories;
\item $F$ is the no-global-frame refusal row, with the same refusal discipline as \autoref{ch:concrete-instances-noglobalsyncboundary-namecert};
\item $I$ is the inter-Hist invariant readback row, recording that the boundary is stable as an inter-Hist packet rather than a single-history rate;
\item $H$ records componentwise $\hsame$ transport for the displayed rows;
\item $Q$ records $\psame$ stability for the probes used by $D,G,R,S,F,I$;
\item $C$ records displayed $\Cont$ routes from the causal route and gate into the local naming row;
\item $P$ is the $\Pkg$ provenance over $\InterHistRateBoundaryUp$, $\CrossHistCausalRateUp$, $\MaxRateReadGateUp$, $\InterHistInvariantUp$, $\NoGlobalSyncBoundaryUp$, $\CrossHistCausalRouteUp$, $\BHist$, $\hsame$, $\psame$, $\Cont$, $\Pkg$, and $\NameCert$;
\item $N$ is the local $\NameCert_{\InterHistRateBoundaryUp}$ row.
\end{enumerate}
The carrier exports a finite maximum-rate boundary for inter-Hist consumers. It exports no host-side speed, global time coordinate, metric tensor, Lorentz frame, selected synchronization map, or detached rate primitive.
\end{definition}

\begin{theorem}[InterHistRateBoundary NameCert obligations]
\label{thm:inter-hist-rate-boundary-namecert-obligations}
For the carrier of \autoref{def:inter-hist-rate-boundary-carrier}, the local $\NameCert_{\InterHistRateBoundaryUp}$ surface consists exactly of seven finite obligations:
\begin{enumerate}
\item carrier admission for $B=(A,D,G,R,S,F,I,H,Q,C,P,N)$;
\item causal-route exactness, requiring $D$ to be read through the displayed $\CrossHistCausalRouteUp$ and $\CrossHistCausalRateUp$ rows;
\item maximum-rate gate exactness, requiring $G$ to bind $R,S,F$ before any consumer reads the rate row;
\item observer-symmetry exactness through $S$ after componentwise transport;
\item no-global-frame refusal through $F$ before any boundary read is public;
\item invariant readback through $I$, keeping the packet on the inter-Hist surface;
\item transport, probe, route, provenance, and local-name stability through $H,Q,C,P,N$.
\end{enumerate}
\end{theorem}
\begin{proof}
Carrier admission is the projection of the finite tuple in \autoref{def:inter-hist-rate-boundary-carrier}. The rows $A,D,G,R,S,F,I,H,Q,C,P,N$ display the observer boundary, causal route, rate gate, maximum-rate row, symmetry row, refusal row, invariant row, transports, probe stability, routes, provenance, and local name.

Apply the causal-rate surface from \autoref{ch:concrete-instances-crosshistcausalrate-namecert}. It supplies the route discipline for $D$ and prevents the boundary from reading a rate before the cross-history causal route is displayed.

Apply the max-rate read gate from \autoref{ch:concrete-instances-maxratereadgate-namecert}. The gate binds $R$ to $S$ and $F$, so the maximum-rate row is not public until both observer symmetry and no-global-frame refusal have been read.

Finally use the invariant and refusal rows. The invariant row keeps the packet on the inter-Hist surface, while the refusal row excludes global frames, clocks, metric tensors, Lorentz frames, synchronized coordinates, and detached host speeds. The remaining rows $H,Q,C,P,N$ transport, route, package, and name the same finite packet.
\end{proof}

\begin{theorem}[InterHistRateBoundary gated consumer]
\label{thm:inter-hist-rate-boundary-gated-consumer}
Every $\CrossHistCausalRateUp$, $\LorentzFrameRateUp$, $\MinkowskiRateGeometryUp$, or inter-Hist locality consumer of an accepted $\InterHistRateBoundaryUp$ packet may read the maximum-rate row $R$ only through the displayed boundary rows
$$
A,D,G,R,S,F,I,H,Q,C,P,N.
$$
In particular, the read of $R$ passes through the causal route $D$, the gate row $G$, the observer-symmetry row $S$, the no-global-frame refusal row $F$, and the invariant row $I$ before it reaches any downstream consumer.
\end{theorem}
\begin{proof}
Project the accepted boundary carrier. The only maximum-rate row is $R$, and the only displayed path to it is through the route and gate rows $D$ and $G$.

Use \autoref{thm:inter-hist-rate-boundary-namecert-obligations}. The obligations make causal-route exactness, gate exactness, symmetry, refusal, and invariant readback part of one local naming surface.

The max-rate gate of \autoref{ch:concrete-instances-maxratereadgate-namecert} supplies the public lock for $R,S,F$. The causal-rate packet of \autoref{ch:concrete-instances-crosshistcausalrate-namecert} supplies the cross-history route read before the rate row is inspected.

Thus a downstream consumer receives a finite boundary packet rather than a free maximum-rate primitive. Any read that bypasses $D$, $G$, $S$, $F$, or $I$ is outside the accepted $\InterHistRateBoundaryUp$ carrier.
\end{proof}

\begin{theorem}[InterHistRateBoundary nonescape]
\label{thm:inter-hist-rate-boundary-nonescape}
An accepted $\InterHistRateBoundaryUp$ packet has no admissible export of a host speed, global time coordinate, synchronized observer frame, metric tensor, Lorentz frame, selected synchronization map, or detached maximum-rate row. Its public surface is exhausted by $A,D,G,R,S,F,I,H,Q,C,P,N$ and the local $\NameCert_{\InterHistRateBoundaryUp}$ obligations over those rows.
\end{theorem}
\begin{proof}
Start from the gated consumer theorem. Every rate read factors through the displayed route, gate, symmetry, refusal, and invariant rows.

The no-global-frame refusal row $F$ is read with the boundary discipline of \autoref{ch:concrete-instances-noglobalsyncboundary-namecert}. It blocks any interpretation of the packet as a global synchronization surface or privileged frame.

The carrier inventory then completes the non-escape argument. Since the tuple contains no coordinate for host speed, global time, metric geometry, Lorentz-frame data, selected synchronization, or detached rate values, no local NameCert projection can export one.
\end{proof}

\begin{theorem}[InterHistRateBoundary NameCert package]
\label{thm:inter-hist-rate-boundary-namecert-package}
The carrier, local NameCert obligations, gated consumer theorem, and non-escape theorem assemble the concrete $\NameCert_{\InterHistRateBoundaryUp}$ package over $\CrossHistCausalRateUp$, $\MaxRateReadGateUp$, $\InterHistInvariantUp$, $\NoGlobalSyncBoundaryUp$, $\CrossHistCausalRouteUp$, $\MaxCausalRateUp$, $\BHist$, $\hsame$, $\psame$, $\Cont$, $\Pkg$, and $\NameCert$. The Lean carrier surface is realized by \texttt{Nontrivial InterHistRateBoundaryUp} and \texttt{FieldFaithful InterHistRateBoundaryUp} instances in \texttt{lean4/BEDC/Derived/InterHistRateBoundaryUp/TasteGate.lean}, with the canonical declaration surface at \texttt{BEDC.Derived.InterHistRateBoundaryUp.taste\_gate}.
\end{theorem}
\begin{proof}
Combine the carrier definition with the NameCert obligations. They fix the twelve displayed rows and bind route, gate, symmetry, refusal, invariant readback, transport, probe stability, continuation, provenance, and local naming.

The gated consumer theorem identifies how downstream rate and geometry consumers may read the packet. The non-escape theorem identifies what they cannot read.

The cited Lean TasteGate carrier supplies the finite event-flow presentation, round-trip readback, nontriviality, and field faithfulness for the same row surface. Therefore the paper package and the Lean carrier use the same finite BEDC data rather than a host metric or global-frame primitive.
\end{proof}
\leanchecked{BEDC.Derived.InterHistRateBoundaryUp.InterHistRateBoundaryTasteGate\_single\_carrier\_alignment}
